% Første linje er dokumentklassen. Brug *altid* memoir!
\documentclass[a4paper,oneside,article]{memoir}

% Herunder indlæses forskellige pakker
\usepackage[utf8]{inputenc}  % Korrekt håndtering af æ, ø og å
\usepackage[T1]{fontenc}  % Korrekt håndtering af æ, ø og å
\usepackage{microtype}  % Typografisk magi! Giver bl.a. pænere orddeling
\usepackage{graphicx}  % Gør det muligt at indsætte billeder
\usepackage{amsmath}  % Giver adgang til uundværlige matematikting
\usepackage{mathtools}  % Retter ting i ams og giver ekstra muligheder
\usepackage[danish]{babel}  % Danske betegnelser og orddeling
\renewcommand{\danishhyphenmins}{22}  % Bedre dansk orddeling
\usepackage[exponent-product=\ensuremath{\cdot}]{siunitx} % Sientific notation for forskellige ting, med en regel om at den skal bruge cdot istedet for times. 
\usepackage{derivative} % Til at skrive differentialligninger 
\usepackage{xcolor} % For at få flere tekst farver 
\usepackage[colorlinks=true, urlcolor=blue]{hyperref} % For at få refrences man kan trykke på, slet [] hvis man vil have kasser i stedet for farvet tekst. 


% Indholdet i dit dokument starter her!
\begin{document}
\author{Mikkel Moth Billing \\ AU737195}
\title{Title}
\date{\today}
\maketitle   % Indsætter overskriften
\hypersetup{linkcolor=black} % Laver tableofcontents sort.  
\tableofcontents*
\hypersetup{linkcolor=red} % og så alle andre hyperrefs røde. 

% Lav ny pagestyle (kun odd fordi enkeltsidet)
\makepagestyle{Title}
\pagestyle{Title}
\makeoddhead{Title}{\color{lightgray}{\theauthor}}{}{\textcolor{lightgray}{\thedate}}
\makeoddfoot{Title}{}{\color{lightgray}{\thepage{} / \thelastpage}}{}
% Også fod på siden med \maketitle
\makeoddfoot{plain}{}{\color{lightgray}{\thepage{} / \thelastpage}}{}
\newpage

%Skriv


\end{document} % ----------Dokumentet er nu slut! ------------